% GENERATED FILE. Edit canonical lessons, then run npm run generate:books.
% canonical-source-hash: fnv1a64:cb7196904c179b89
% canonical-lessons: TE-C360-khadgam, TE-C360-kavacam, TE-C360-ayudham, TE-C360-villu, TE-C360-tupaki

\chapter{Weapons}
\label{ch:te-weapons}
\hlchaptermodality{360}

\begin{chapteropening}
\textbf{By the end of this chapter:} \emph{I can say a sword, armour, a weapon, a bow and a gun.}

Complete the last lesson of chapter 360: I can say a sword, armour, a weapon, a bow and a gun.
\end{chapteropening}

\section[\texorpdfstring{\te{ఖడ్గం} (khaḍgaṁ)}{ఖడ్గం (khaḍgaṁ)}]{\te{ఖడ్గం} (khaḍgaṁ) — a sword}

\label{lesson:TE-C360-khadgam}



\begin{quote}
Before the new one: say the Telugu for a princess, then the Telugu for to found.
\end{quote}

\subsection*{You'll want to know: \te{ఖడ్గం}}
\textbf{\te{ఖడ్గం}} — \emph{khaḍgaṁ} — \textquotedblleft{}a sword\textquotedblright{}.

\begin{quote}
\te{సైనికుడు} \te{ఖడ్గం} \te{పట్టుకుని} \te{నిలబడ్డాడు}.
\end{quote}

\begin{quote}
\emph{sainikuḍu khaḍgaṁ paṭṭukuni nilabaḍḍāḍu.}
\end{quote}

\begin{quote}
The soldier stood holding a sword.
\end{quote}

A word of stories; at home a knife is \te{కత్తి}.

\begin{cousinweb}
\textbf{\te{ఖడ్గం}} (\emph{khaḍgaṁ}) is borrowed from Sanskrit \emph{khaḍga}, \textquotedblleft{}a sword\textquotedblright{} (CDIAL 3787).
\end{cousinweb}

\begin{grammarlens}[title={What you've built}]
One more word for telling what happened.
\end{grammarlens}

\subsection*{Your turn}
\begin{itemize}
\raggedright
  \item \emph{Say it:} \emph{khaḍgaṁ}
  \item \emph{Say it:} \emph{khaḍgaṁ}, once more
  \item \emph{Say it:} \emph{sainikuḍu khaḍgaṁ paṭṭukuni nilabaḍḍāḍu.} — The soldier stood holding a sword.
  \item \emph{Recall:} say the Telugu for a princess, then the Telugu for to found, then say \emph{khaḍgaṁ} again
\end{itemize}

\begin{tcolorbox}[breakable,colback=teal!4,colframe=teal!35!black,title={Before you move on}]
What does \te{ఖడ్గం} mean? (\textquotedblleft{}A sword\textquotedblright{}.) Say it once more.
\end{tcolorbox}
\section[\texorpdfstring{\te{కవచం} (kavacaṁ)}{కవచం (kavacaṁ)}]{\te{కవచం} (kavacaṁ) — armour; a shield}

\label{lesson:TE-C360-kavacam}



\begin{quote}
Before the new one: say the Telugu for to found, then the Telugu for a sword.
\end{quote}

\subsection*{You'll want to know: \te{కవచం}}
\textbf{\te{కవచం}} — \emph{kavacaṁ} — \textquotedblleft{}armour; a shield\textquotedblright{}.

\begin{quote}
\te{సైనికులు} \te{ఇనుప} \te{కవచాలు} \te{ధరించారు}.
\end{quote}

\begin{quote}
\emph{sainikulu inupa kavacālu dhariñcāru.}
\end{quote}

\begin{quote}
The soldiers wore iron armour.
\end{quote}

\begin{cousinweb}
\textbf{\te{కవచం}} (\emph{kavacaṁ}) is borrowed from Sanskrit \emph{kavaca}, \textquotedblleft{}a coat of mail\textquotedblright{} (CDIAL 2957); CDIAL notes a view that the Sanskrit word itself came from Dravidian.
\end{cousinweb}

\begin{grammarlens}[title={What you've built}]
One more word for telling what happened.
\end{grammarlens}

\subsection*{Your turn}
\begin{itemize}
\raggedright
  \item \emph{Say it:} \emph{kavacaṁ}
  \item \emph{Say it:} \emph{kavacaṁ}, once more
  \item \emph{Say it:} \emph{sainikulu inupa kavacālu dhariñcāru.} — The soldiers wore iron armour.
  \item \emph{Recall:} say the Telugu for to found, then the Telugu for a sword, then say \emph{kavacaṁ} again
\end{itemize}

\begin{tcolorbox}[breakable,colback=teal!4,colframe=teal!35!black,title={Before you move on}]
What does \te{కవచం} mean? (\textquotedblleft{}Armour; a shield\textquotedblright{}.) Say it once more.
\end{tcolorbox}
\section[\texorpdfstring{\te{ఆయుధం} (āyudhaṁ)}{ఆయుధం (āyudhaṁ)}]{\te{ఆయుధం} (āyudhaṁ) — a weapon}

\label{lesson:TE-C360-ayudham}



\begin{quote}
Before the new one: say the Telugu for a sword, then the Telugu for armour.
\end{quote}

\subsection*{You'll want to know: \te{ఆయుధం}}
\textbf{\te{ఆయుధం}} — \emph{āyudhaṁ} — \textquotedblleft{}a weapon\textquotedblright{}.

\begin{quote}
\te{దొంగల} \te{దగ్గర} \te{ఆయుధాలు} \te{దొరికాయి}.
\end{quote}

\begin{quote}
\emph{doṅgala daggara āyudhālu dorikāyi.}
\end{quote}

\begin{quote}
Weapons were found on the robbers.
\end{quote}

Plural: \te{ఆయుధాలు}, \emph{āyudhālu}.

\begin{cousinweb}
\textbf{\te{ఆయుధం}} (\emph{āyudhaṁ}) is borrowed from Sanskrit \emph{āyudha}, \textquotedblleft{}a weapon\textquotedblright{} (CDIAL 1291), from \emph{yudh}, \textquotedblleft{}to fight\textquotedblright{}, the root of \te{యుద్ధం}.
\end{cousinweb}

\begin{grammarlens}[title={What you've built}]
One more word for telling what happened.
\end{grammarlens}

\subsection*{Your turn}
\begin{itemize}
\raggedright
  \item \emph{Say it:} \emph{āyudhaṁ}
  \item \emph{Say it:} \emph{āyudhaṁ}, once more
  \item \emph{Say it:} \emph{doṅgala daggara āyudhālu dorikāyi.} — Weapons were found on the robbers.
  \item \emph{Recall:} say the Telugu for a sword, then the Telugu for armour, then say \emph{āyudhaṁ} again
\end{itemize}

\begin{tcolorbox}[breakable,colback=teal!4,colframe=teal!35!black,title={Before you move on}]
What does \te{ఆయుధం} mean? (\textquotedblleft{}A weapon\textquotedblright{}.) Say it once more.
\end{tcolorbox}
\section[\texorpdfstring{\te{విల్లు} (villu)}{విల్లు (villu)}]{\te{విల్లు} (villu) — a bow}

\label{lesson:TE-C360-villu}



\begin{quote}
Before the new one: say the Telugu for armour, then the Telugu for a weapon.
\end{quote}

\subsection*{You'll want to know: \te{విల్లు}}
\textbf{\te{విల్లు}} — \emph{villu} — \textquotedblleft{}a bow\textquotedblright{}.

\begin{quote}
\te{వేటగాడి} \te{దగ్గర} \te{ఒక} \te{విల్లు} \te{ఉంది}.
\end{quote}

\begin{quote}
\emph{vēṭagāḍi daggara oka villu undi.}
\end{quote}

\begin{quote}
The hunter has a bow.
\end{quote}

A hunter's weapon in the old stories.

\begin{cousinweb}
\textbf{\te{విల్లు}} (\emph{villu}) is native Dravidian: DEDR 5422 lists Telugu \emph{villu}, \textquotedblleft{}a bow\textquotedblright{}, with Tamil \emph{vil}.
\end{cousinweb}

\begin{grammarlens}[title={What you've built}]
One more word for telling what happened.
\end{grammarlens}

\subsection*{Your turn}
\begin{itemize}
\raggedright
  \item \emph{Say it:} \emph{villu}
  \item \emph{Say it:} \emph{villu}, once more
  \item \emph{Say it:} \emph{vēṭagāḍi daggara oka villu undi.} — The hunter has a bow.
  \item \emph{Recall:} say the Telugu for armour, then the Telugu for a weapon, then say \emph{villu} again
\end{itemize}

\begin{tcolorbox}[breakable,colback=teal!4,colframe=teal!35!black,title={Before you move on}]
What does \te{విల్లు} mean? (\textquotedblleft{}A bow\textquotedblright{}.) Say it once more.
\end{tcolorbox}
\section[\texorpdfstring{\te{తుపాకి} (tupāki)}{తుపాకి (tupāki)}]{\te{తుపాకి} (tupāki) — a gun}

\label{lesson:TE-C360-tupaki}



\begin{quote}
Before the new one: say the Telugu for a weapon, then the Telugu for a bow.
\end{quote}

\subsection*{You'll want to know: \te{తుపాకి}}
\textbf{\te{తుపాకి}} — \emph{tupāki} — \textquotedblleft{}a gun\textquotedblright{}.

\begin{quote}
\te{పోలీసు} \te{దగ్గర} \te{తుపాకి} \te{ఉంది}.
\end{quote}

\begin{quote}
\emph{pōlīsu daggara tupāki undi.}
\end{quote}

\begin{quote}
The policeman has a gun.
\end{quote}

Plural: \te{తుపాకులు}, \emph{tupākulu}.

\te{తుపాకి} \te{పేల్చు}, to fire a gun.

\begin{cousinweb}
\textbf{\te{తుపాకి}} (\emph{tupāki}) is the Telugu word for a gun; its origin is not settled.
\end{cousinweb}

\begin{grammarlens}[title={What you've built}]
One more word for telling what happened.
\end{grammarlens}

\subsection*{Your turn}
\begin{itemize}
\raggedright
  \item \emph{Say it:} \emph{tupāki}
  \item \emph{Say it:} \emph{tupāki}, once more
  \item \emph{Say it:} \emph{pōlīsu daggara tupāki undi.} — The policeman has a gun.
  \item \emph{Recall:} say the Telugu for a weapon, then the Telugu for a bow, then say \emph{tupāki} again
\end{itemize}

\begin{tcolorbox}[breakable,colback=teal!4,colframe=teal!35!black,title={Before you move on}]
What does \te{తుపాకి} mean? (\textquotedblleft{}A gun\textquotedblright{}.) Say it once more.
\end{tcolorbox}
